\section*{Anexo A. Datos de la instancia}
\addcontentsline{toc}{section}{Anexo A. Datos de la instancia}

\subsection*{A.1 Actividades}
Esta tabla describe las actividades candidatas, sus características, costes,
obligatoriedad, espacios compatibles y franjas disponibles.
\begin{table}[htbp]
\centering
\scriptsize
\resizebox{\textwidth}{!}{%
\begin{tabular}{rlccrrrrrl}
\toprule
Id & Actividad & Dur. & Categoría & Asist. & Valor. & Coste & Obl. & Esp. & Franjas \\
\midrule
1 & Conferencia inaugural: IA y sociedad & 1 & ML & 180 & 4.2 & 1200 & Sí & E1 & 1--2 \\
2 & Conferencia de clausura: datos y ética & 1 & Estadística & 160 & 4.0 & 1000 & Sí & E1 & 7--8 \\
3 & Taller de redes neuronales & 3 & ML & 35 & 4.8 & 900 & No & E2 & 1,2,3,5--8 \\
4 & Taller de pandas y limpieza de datos & 2 & Ingeniería & 38 & 4.2 & 500 & No & E2 & Todas \\
5 & Hackathon de datos & 3 & Ingeniería & 40 & 4.9 & 1100 & No & E2 & 5--8 \\
6 & Taller de inferencia bayesiana & 2 & Estadística & 30 & 4.7 & 550 & No & E2 & Todas \\
7 & Charla: IA generativa & 1 & ML & 190 & 3.6 & 1800 & No & E1 & 7 \\
8 & Charla: pipelines de datos en producción & 1 & Ingeniería & 45 & 3.0 & 350 & No & E1,E3 & Todas \\
9 & Charla: series temporales & 1 & Estadística & 40 & 3.0 & 250 & No & E1,E3 & 4 \\
10 & Mesa redonda: datos y empleo & 2 & Ingeniería & 110 & 3.4 & 900 & No & E1 & 2--4 \\
11 & Charla: visualización de datos & 1 & Estadística & 35 & 2.8 & 250 & No & E1,E3 & Todas \\
12 & Charla: aprendizaje por refuerzo & 1 & ML & 50 & 3.2 & 350 & No & E1,E3 & 6--8 \\
13 & Taller de SQL avanzado & 2 & Ingeniería & 30 & 4.0 & 400 & No & E2 & Todas \\
14 & Charla: Big Data con Spark & 1 & Ingeniería & 60 & 3.5 & 450 & No & E1,E3 & 1--4 \\
15 & Charla: contrastes y A/B testing & 1 & Estadística & 30 & 2.7 & 200 & No & E1,E3 & Todas \\
16 & Charla: explicabilidad en ML & 1 & ML & 85 & 3.6 & 500 & No & E1 & 6--8 \\
17 & Seminario de estadística espacial & 2 & Estadística & 45 & 3.0 & 350 & No & E1,E3 & 3--4 \\
18 & Charla: ética y sesgos algorítmicos & 1 & ML & 75 & 4.0 & 1000 & No & E1,E3 & 7 \\
\bottomrule
\end{tabular}%
}
\caption{Actividades de la instancia principal.}
\end{table}

\subsection*{A.2 Espacios y compatibilidades}
Esta tabla recoge los espacios disponibles, su capacidad y las franjas en las
que permanecen abiertos.
\begin{table}[htbp]
\centering
\begin{tabular}{llrl}
\toprule
Id & Espacio & Capacidad & Franjas abiertas \\
\midrule
E1 & Auditorio & 200 & 1--8 \\
E2 & Aula de ordenadores & 40 & 1,2,3,5--8 \\
E3 & Sala de seminarios & 80 & 2,3,4,6--8 \\
\bottomrule
\end{tabular}
\caption{Espacios de la instancia principal.}
\end{table}

Las compatibilidades se almacenan en \texttt{compat\_espacios.csv}. Además,
la compatibilidad de tipo debe combinarse con la capacidad: las actividades
7, 10 y 16 son compatibles con E3 por tipo, pero su asistencia supera las
80 plazas y no pueden asignarse a ese espacio.

\subsection*{A.3 Incompatibilidades y parámetros}
Esta tabla resume las parejas de actividades incompatibles y los parámetros
globales utilizados por el modelo.
\begin{table}[htbp]
\centering
\begin{tabular}{ll}
\toprule
Tipo & Valor \\
\midrule
Incompatibilidades & (3,10), (12,16), (9,17), (7,18), (8,14) \\
Presupuesto & 6500 \\
Pesos por defecto & $w_A=0.4$, $w_V=0.4$, $w_N=0.2$ \\
Franjas & 8; 4 franjas por día \\
\bottomrule
\end{tabular}
\caption{Parámetros globales e incompatibilidades.}
\end{table}

\clearpage
\section*{Anexo B. Contribución de cada integrante}
\addcontentsline{toc}{section}{Anexo B. Contribución de cada integrante}

La siguiente tabla resume las tareas realizadas por cada integrante a partir
del plan de trabajo de cada parte.

\begin{table}[htbp]
\centering
\renewcommand{\arraystretch}{1.2}
\begin{tabular}{p{0.22\textwidth}p{0.12\textwidth}p{0.56\textwidth}}
\toprule
\textbf{Integrante} & \textbf{Parte} & \textbf{Contribución} \\
\midrule
Héctor Nieto Colás & 1 &
\textbf{Objetivo:} definir el contexto, los datos y el modelo matemático que
selecciona, asigna y programa las actividades del evento.
\newline
\textbf{Asignaciones:} elegir y describir el contexto; presentar franjas,
espacios, actividades y restricciones naturales; caracterizar actividades y
espacios; aportar los valores y documentar los CSV externos; incluir las tablas
completas en el anexo; definir y justificar los tres criterios y sus pesos;
incorporar y justificar las incompatibilidades temporales; y definir conjuntos,
parámetros, variables, función objetivo y restricciones R1--R9. \\
\midrule
Nicolás Mínguez Aguirre & 2 &
\textbf{Objetivo:} convertir el modelo formulado en una implementación PLE
reproducible y una interfaz Streamlit mínima.
\newline
\textbf{Asignaciones:} cargar y validar los siete CSV; construir y resolver el
modelo binario con PuLP e HiGHS; devolver selección, asignación e indicadores;
distinguir solución óptima, inviabilidad y errores; implementar la interfaz,
sus controles, calendario, tablas y descarga CSV; incorporar la captura; y
documentar la ejecución de Streamlit y la compilación del informe. \\
\midrule
Unai Lara Domínguez & 3 &
\textbf{Objetivo:} redactar el análisis de los resultados y documentar el
software utilizado.
\newline
\textbf{Asignaciones:} redactar las expectativas de cada escenario; ejecutar
el modelo y tabular los indicadores; comentar los compromisos entre criterios;
documentar la reproducibilidad; verificar la gestión de inviabilidad, la
exportación CSV y la coincidencia entre interfaz e informe; incorporar la
captura; sintetizar el valor del modelo y los compromisos; y redactar la
declaración de inteligencia artificial y autoría humana. \\
\bottomrule
\end{tabular}
\caption{Contribución individual según los planes de trabajo de las tres partes.}
\end{table}
